\section{Method}
\subsection{Data and cohort construction}
Source A is EBBU-2017 (\res{nAn} URLs: \res{nAbenign} legitimate, \res{nAphishing} phishing; repository \texttt{ebubekirbbr/pdd} at commit \texttt{90b394d}); the class is given by the file name. Source B is HP (\res{nBn} URLs), rebuilt by concatenating the \texttt{hp\_train/valid/test} files of \texttt{fariza25/PhishingURLDetection} at commit \texttt{ff9f1f5}; the original split is kept only as provenance. Repository B has no numeric codebook. Its copy of EBBU agrees with source A on \res{labelJoin} joined URLs with \res{labelAgree}\% agreement, supporting the repository convention 1 = phishing. An additional evaluator-only audit uniquely matches both numeric classes to the benign/phishing counts in every original HP split in Table~1 of~\cite{rashid2024}, supporting 0 = benign and 1 = phishing directly against published class aggregates. This is count-based crossvalidation; an author-written numeric codebook is not supplied. All file hashes, sizes and observed counts are stored in a manifest; observed counts equal the expected ones.

URLs are canonicalised offline (no network access is possible during extraction): outer whitespace is trimmed; strings with control characters, schemes other than HTTP(S), invalid ports or unparsable hosts are quarantined (\res{quarA} in A, \res{quarB} in B); hosts are lower-cased and IDNA-encoded, while path, query and fragment keep their case and percent-encoding. Exact duplicates (same canonical key and label) are removed within each source (\res{dupA} in A, \res{dupB} in B), and keys with conflicting labels would be quarantined in both sources (none occurred). Each URL is grouped by its registered domain under the ICANN section of a frozen Public Suffix List snapshot; private suffixes are deliberately ignored so that, e.g., all \texttt{blogspot.com} sites form one conservative group. The two sources shared \res{sharedGroups} registered domains (only \res{exactOverlap} exact URLs); all URLs of these domains were removed from \emph{both} sources (\res{purgeA} from A, \res{purgeB} from B). The purge consults target identities only for this audit, never for learning, and costs mostly benign URLs (\res{lossBenignA}\% of A's and \res{lossBenignB}\% of B's benign URLs, versus \res{lossPhishA}\% and \res{lossPhishB}\% of phishing), because popular legitimate platforms appear in both. Table~\ref{tab:attrition} gives the full attrition. The resulting cohorts have \res{cAn} (A) and \res{cBn} (B) URLs over \res{groupsA} and \res{groupsB} domains; the largest group is \texttt{\res{topGroupA}} (\res{topGroupShareA}\%) in A and \texttt{\res{topGroupB}} (\res{topGroupShareB}\%) in B.

\begin{table}[t]
\caption{Cohort attrition (each URL counted once, at the first stage that removed it).}
\label{tab:attrition}
\centering\scriptsize
\begin{tabular}{llrrr}
\toprule
Source & Stage & URLs & Benign & Phishing \\
\midrule
EBBU-2017 & Loaded & 73,575 & 36,400 & 37,175 \\
EBBU-2017 & Parsed & 73,572 & 36,400 & 37,172 \\
EBBU-2017 & Deduplicated & 54,761 & 29,910 & 24,851 \\
EBBU-2017 & Conflicts removed & 54,761 & 29,910 & 24,851 \\
EBBU-2017 & Cohort (shared domains purged) & 46,769 & 22,167 & 24,602 \\
\midrule
HP & Loaded & 92,500 & 46,250 & 46,250 \\
HP & Parsed & 92,483 & 46,233 & 46,250 \\
HP & Deduplicated & 92,482 & 46,233 & 46,249 \\
HP & Conflicts removed & 92,482 & 46,233 & 46,249 \\
HP & Cohort (shared domains purged) & 75,077 & 29,758 & 45,319 \\
\bottomrule
\end{tabular}

\end{table}

\subsection{Splits}
Each source is split into FIT/VAL/TEST ($\approx$60/20/20\% of URLs) by a deterministic greedy stratified \emph{group} assignment (seed 20261115): whole domains are placed, largest first, into the partition whose per-class quota is most unfilled. Partitions share no domain, within or across sources, and keep the class prevalence (Table~\ref{tab:splits}). FIT is used for selection and fitting, VAL only for thresholds, and TEST once, after freezing.

\begin{table}[t]
\caption{Group-disjoint partitions.}
\label{tab:splits}
\centering\scriptsize
\begin{tabular}{llrrrr}
\toprule
Source & Part & URLs & Domains & \% URLs & \% phish \\
\midrule
EBBU & FIT & 28,061 & 8,018 & 60.0 & 52.6 \\
EBBU & VAL & 9,355 & 6,448 & 20.0 & 52.6 \\
EBBU & TEST & 9,353 & 6,447 & 20.0 & 52.6 \\
HP & FIT & 45,050 & 2,614 & 60.0 & 60.4 \\
HP & VAL & 15,014 & 2,060 & 20.0 & 60.4 \\
HP & TEST & 15,013 & 2,061 & 20.0 & 60.4 \\
\bottomrule
\end{tabular}

\end{table}

\subsection{Features}
We use 37 deterministic lexical features (lengths of URL, host, path, query and fragment; counts of 10 delimiter characters, digits, letters and non-ASCII code points; digit ratios; character entropy; longest alphanumeric run; host label and subdomain counts; IP host, explicit port and port number; user-info, query, HTTPS and missing-scheme flags; punycode labels; percent-encoded triplets; query pairs; path segments; and token counts of a fixed 10-word keyword list in the URL and in the host). Raw strings, domains, dataset identity and split labels never enter a model; no vocabulary or target encoding is learned. Each feature is a separate function, so a subset computes only what it needs, while parsing is always paid.

\subsection{Source-only feature selection}
All selectors use the same RF proxy (200 trees, depth 16, min.\ leaf 2, $\sqrt{p}$ features). FIT is divided into three stratified \emph{domain} folds; in each fold the proxy is trained on two folds and the permutation importance (PI) of every feature is measured on the third as the drop in average precision over five permutations. MEAN-PI ranks features by mean PI over the folds. STABLE-PI first counts, for each feature, the folds in which it is in the top $q{=}12$ \emph{and} has PI${>}0$ (votes), then ranks by votes, median PI, the median absolute deviation of its fold ranks, and feature ID. The budget $k\in\{8,16,24\}$ is chosen by nested cross-validation: in three outer domain folds of FIT, PI is recomputed from the outer-training part only, each classifier is trained on the top-$k$ features and scored on the outer hold-out; the smallest $k$ within 0.01 AP of the best mean is kept, separately for every source, classifier, selector and seed. PI is then recomputed on all of FIT for the final subset. To isolate the effect of stability from that of the budget, every STABLE-PI model is paired with a MEAN-PI model at the \emph{same} $k$ (matched-$k$ contrast).

\subsection{Classifiers, thresholds and evaluation}
RF uses 300 trees, depth 20, min.\ leaf 2, $\sqrt{p}$ features; LightGBM uses 300 rounds, learning rate 0.05, 31 leaves, no subsampling, $\lambda_2{=}1$ and deterministic CPU mode. Hyper-parameters are fixed a priori and are not claimed to be optimal. Models are trained on source FIT with seeds 17, 29 and 43. On source VAL we lock (i) the threshold maximising F1 (primary), (ii) a threshold maximising recall subject to VAL false-positive rate $\le$1\%, with ties resolved by the largest threshold, and (iii) 0.5 (sensitivity). The same model and threshold are then applied to the source TEST (internal) and to the other source's TEST (cross), with no retraining or recalibration on the target. The transfer drop is F1(source TEST)$-$F1(target TEST).

The pre-specified primary outcome is target-TEST F1 of STABLE-PI minus MEAN-PI at matched $k$, for each classifier and direction (four contrasts). Uncertainty is a paired domain-cluster bootstrap (\res{nBoot} replicates): target domains are resampled with replacement, both arms and all three seeds use the same resample, and the per-replicate delta is averaged over seeds. These intervals are conditional on the fixed split and seeds, and we make no family-wise significance claim; an absolute F1 difference of 0.01 was fixed in advance as the smallest difference worth discussing.

\subsection{Cost and attributions}
Latency is measured on the seed-17 model of every configuration with the same interleaved VAL URLs of both sources: 10 warm-up and 300 (batch 1) or 30 (batch 1024) timed repetitions, reporting medians for parsing+feature extraction, prediction alone and end-to-end, with one and with four inference threads. SHAP values are computed with TreeExplainer on the actual classifiers (interventional, probability output, 100 source-FIT background rows) for 500 label-stratified source-TEST and 500 target-TEST URLs, for the seed-17 ALL and STABLE-PI models. The explanations here are an explicitly exploratory post-freeze correction on exactly the saved feature matrices. RF inputs/background are cast to the classifier's float32 input lattice; SHAP's threshold representation is rounded downward on that lattice to preserve the original classifier's split decisions. The fitted classifiers are unchanged. Every corrected cohort passes the existing $10^{-6}$ tolerance, with maximum residual \res{shapAddErr}; explainer predictions and background expectations are independently checked against the actual classifier. Additivity is a numerical consistency property and does not by itself establish explanation fidelity.

\subsection{Baselines and exploratory diagnostics (RQ5)}
After the frozen evaluation we added diagnostics that were specified before they were run and that use \emph{only} FIT and VAL; no TEST feature or label is read, and the frozen models are loaded with hash verification. \emph{Baselines:} (a) L2-regularised logistic regression (LogReg) on all 37 features after $\log(1{+}x)$ and standardisation, $C{=}1$, without tuning; (b) a one-feature rule that flags every non-HTTPS URL. \emph{Shortcut ablation:} RF and LightGBM are refitted on source FIT without F30 (\texttt{is\_https}), with the frozen hyper-parameters and seeds. LogReg and the ablated models take their F1-max threshold from source VAL, like the frozen models. Transfer is measured on the \emph{target VAL} partition, which no model of that direction has seen; source-VAL F1 is optimistic because the same data set the threshold. Differences use the paired domain-cluster bootstrap above. \emph{Feature analysis:} for 14 interpretable features we compute Cohen's $d$ (phishing minus benign) on each source's FIT. \emph{Error analysis:} for the seed-17 frozen ALL models we report false-positive and false-negative rates in pre-defined slices (HTTPS, URL length, subdomain count, query) and the five most confident errors of each type. URLs are shown defanged and are never opened. These results are exploratory and conditional on one split.

\subsection{DA-GAFS: direction- and generalization-aware selection (RQ6)}\label{sec:dagafs}
DA-GAFS (Direction-Aware, Generalization-Aware Feature Selection) organises the selection problem into five components, specified before running and evaluated only on FIT/VAL. Unlike the frozen study, it assumes that a \emph{second labelled collection} is available for choosing features, but not for training: selection uses the labelled FIT partitions of both collections, every model is fitted on source FIT only, thresholds are locked on source VAL, and transfer is measured on the domain-disjoint target VAL. \emph{C1, cross-source stability.} From the saved full-FIT PI vectors of both sources (3 seeds $\times$ 3 folds each, hash-verified against the freeze), $S_f$ is the share of the 18 vectors in which feature $f$ is in the top $q{=}12$ with PI${>}0$; pooled importance $I_f$ averages the max-normalised mean PI of the two sources. \emph{C2, directional consistency.} On each source's FIT we compute the rank-biserial association $r_{s,f}=2\,\mathrm{AUC}_{s,f}-1$ between the feature and the phishing label; a feature conflicts if $r_{A,f}r_{B,f}<0$ and $\min_s|r_{s,f}|\ge0.10$, and then $D_f{=}0$ (else 1). Sign stability is checked with 200 domain-bootstrap replicates. \emph{C3, unified score.} $G_f=D_f\,(I_f+S_f)/2$; ablations rank by $I_f$ alone (XS-I) and by $(I_f+S_f)/2$ without the direction gate (XS-IS). \emph{C4, multi-objective budget.} For every prefix $k=1,\dots,37$ of each ranking, the generalisation objective is the mean over both directions of the AP of a LightGBM (v1 parameters, seed 17) trained on source FIT and scored on target FIT; the cost objective is the measured parsing time plus the sum of per-feature extraction times ($\mu$s per URL on 4{,}000 FIT URLs). The chosen $k^*$ is the cheapest prefix within 0.01 AP of that ranking's best, the same tolerance as the frozen budget rule. \emph{C5, generalization-aware ensemble.} RF, LightGBM and the LogReg recipe of RQ5 are trained on the DA-GAFS features; every non-empty combination is scored by equal-weight probability averaging on target FIT, its cost adds the measured prediction time of its members, and the cheapest combination within 0.01 AP of the best is kept. Paired contrasts on target VAL (DA-GAFS versus ALL and STABLE-PI of the frozen study, and the component ablations) use the bootstrap above. After seeing the C4 output we additionally evaluated, as a labelled post-hoc analysis, the prefix that the same rule chose for each ablation ranking.

\subsection{Protocol control}
Code, configurations, data splits, feature matrices and all fitted models of both directions were hashed into a freeze manifest (hash \texttt{\res{freezeHash}}) before the evaluator read any TEST label; the development loader cannot load TEST at all. Reported values are generated from saved evaluation/analysis artefacts. The original I06 check recomputed detection metrics; later audits independently reconciled aggregate counts, model probabilities and presentation lineage without rereading TEST labels. This is a procedural control rather than a cryptographic guarantee.
